% GENERATED by paper/integrate_figures_v2.py from figures/v7/supp/caption_supp.tex (v2 figure caption, rendered by the figure script).
% Do not edit: change the figure script and re-run build.sh.
\caption{\textbf{Shopify's Web Pixels Manager (WPM) decides in the page's main thread whether a sandboxed pixel starts, and never stops it afterwards.} A schematic of our static reading of the WPM manager bundle and its inline loader (Methods), read from left to right and drawn like Fig.~1a: the outer frame is the Shopify store in the visitor's browser; inside it, the light grey ground is the page's main thread and the darker grey ground beside it the pixel sandboxes, which are isolated from the page; the outlined frame inside the page is the manager bundle. Boxes are components, with their secondary text in grey; arrows show what feeds what. Inline loader (top, in the page): each storefront page lists every pixel (`pixel list') with its type (`APP' or `CUSTOM'), its runtime context (`runtimeContext': `STRICT', `LAX' or `OPEN') and its declared privacy purposes (analytics, marketing, preferences, sale of data), so a crawler can read the list before any script runs; the list feeds the consent gate. Consent signals (left), in the order and words of Fig.~1a: the choice of a consent-management platform (CMP) that calls `setTrackingConsent'; the merchant's region default, written by the server into the page response; and, for sale of data only, Global Privacy Control (GPC), whose `Sec-GPC' request header makes the server decline sale of data. Manager bundle (middle): the consent state is the first one found of the sources listed in its box, in this order: `cachedConsent', `injectedConsent', the `\_cs' URL parameter, the `\_tracking\_consent' cookie and the `Server-Timing: \_cmp' header; if none can be read, everything is allowed. The event bus replays page events to every pixel, and each event carries the full page URL. The consent gate starts a pixel only when the consent state allows every purpose the pixel declares, and starts a pixel that declares no purpose at once. Runtime containers (right; the grey qualifier in each box is the runtime context): the gate starts each pixel in the container its runtime context names, a web worker (strict, `STRICT') or a sandboxed iframe (lax, `LAX'), both on the pixel sandboxes' ground, so a script blocker in the page cannot see them, or a script in the page (open, `OPEN'), drawn on the page's ground, which a blocker can see. In the web worker the vendor library sends with `fetch' without cookies, and `navigator' is overwritten, so pixel code reads GPC as undefined. We found no code path that stops a pixel once it has started. Schematic drawn in TikZ, with no generated imagery.
Drawn by figure-generation scripts written with AI assistance (Section~\ref*{M:sec:method-ai} of the main article).}
\label{fig:s1}
